\section{Introduction}
\label{sec:intro}

A formation of unmanned aerial vehicles holds its shape by exchanging
state. How often that exchange should happen is not a detail: it sets the
radio duty cycle, the energy budget, the number of vehicles a channel can
carry, and --- through information staleness --- the accuracy of the
formation itself.

Two answers dominate practice, and each fails in a way the other does not.

\textbf{Fixed-rate periodic communication} is simple, predictable and easy
to schedule. Its weakness is that the rate must be chosen in advance, for a
network whose quality is generally unknown and often changing. Choose a
low rate and the formation degrades exactly when the channel degrades;
choose a high rate and the radio is busy paying for margin it does not
need. In our own frozen results this is visible directly: a $10$~Hz
periodic baseline transmits \MetricDataPtenClean{}~Hz in a clean network
and the same \MetricDataPtenStressed{}~Hz in a heavily impaired one, while
its formation error grows from \MetricRmsePtenClean{}~m to
\MetricRmsePtenStressed{}~m. The policy does not notice that anything has
changed.

\textbf{Conventional event-triggered communication} transmits when the
local state has moved far enough to be worth reporting. This does save
traffic, and substantially: our state-event baseline runs at
\MetricDataEventStressed{}~Hz in the Stressed condition against
\MetricDataPtwentyStressed{}~Hz for $20$~Hz periodic. But a pure state
trigger is blind to what the receiver actually holds. If a packet is lost
or badly delayed, the sender's own state may be perfectly steady --- no
trigger --- while the receiver's copy grows arbitrarily stale. The
resulting error is severe: \MetricRmseEventStressed{}~m, worse than either
periodic baseline.

The missing quantity is \emph{age of information}: not how much the
sender's state has changed, but how old the receiver's copy of it is. Using
it, however, raises a problem that is easy to state and easy to get wrong.
\textbf{A sender cannot observe the age of information at a receiver.} It
can only infer it, and only from something that comes back.

\subsection*{What this paper contributes}

We present a \emph{causal ACK-assisted, AoI-aware event-triggered
communication policy} for multi-UAV formation control. The name is
deliberately narrow. The policy is decentralised in that each sender
decides alone, using only locally held quantities; it is \emph{not}
feedback-free, and we do not describe it as such. Its estimate of receiver
freshness comes entirely from cumulative acknowledgements that traverse a
reverse channel with its own delay, jitter and loss. There is no
same-timestep read of receiver state anywhere in the implementation, and
nine protocol invariants are checked at runtime to keep it that way.

Three design elements carry the result:

\begin{enumerate}
\item \textbf{Dual sender memory.} The sender keeps the last state it
      \emph{transmitted}, which defines innovation, separately from the
      last state an acknowledgement has \emph{confirmed}, which defines
      estimated freshness. Conflating the two is not a shortcut; it makes
      the policy re-send state that is already in flight.
\item \textbf{Freshness-adaptive thresholds.} As the estimated age of the
      receiver's copy grows past a threshold, the state-change threshold
      required to transmit is reduced continuously toward a floor. Age
      alone never triggers a transmission; it only makes the sender easier
      to trigger.
\item \textbf{Separation of new information from refresh.} A repetition
      cooldown governs only repetition. Traffic that carries genuinely new
      information is not subject to it, and refresh traffic is additionally
      forbidden while any packet remains unacknowledged, so the policy does
      not add a duplicate to something already on the wire.
\end{enumerate}

The third element is the one the design history turns on, and we present
that history rather than hiding it. Two earlier versions of this policy
failed for opposite reasons --- one over-transmitted until it breached its
resource ceiling, the other over-suppressed until it lost adaptivity
entirely --- and both failures are reported in
Section~\ref{sec:results} as an ablation chain measured on frozen data.
Neither was repaired by tuning a parameter. No parameter value differs
between the failing and working variants.

\subsection*{What we claim, and what we do not}

The campaign behind this paper was pre-registered: acceptance gates,
selected operating points, statistical procedure and seed policy were
fixed in writing before each experiment ran, and the final validation used
\MetricSeeds{} seeds never used during development. That discipline makes
the negative results as reportable as the positive ones, and we report
them.

\textbf{We claim} that one fixed configuration adapts its communication
effort across network conditions
(\MetricDataCausalClean{} $<$ \MetricDataCausalModerate{} $<$
\MetricDataCausalStressed{}~Hz), that it improves accuracy over
conventional state-event communication in every cell of the final matrix
(\MetricBeatEventCells{}/\MetricTotalCells{}), that it improves on $10$~Hz
periodic in the Stressed condition with a paired confidence interval
excluding zero, and that in the Moderate condition it occupies a
favourable accuracy-versus-communication trade-off, non-dominated in
\MetricParetoModerateWtwentyfive{}\% of the selected point families under
the reference cost model.

\textbf{We do not claim} universal superiority over optimally tuned
periodic communication. A $20$~Hz periodic baseline attains lower formation
error than our policy in \MetricTotalCells{}$-$\MetricBeatPtwentyCells{}
of \MetricTotalCells{} cells, and our policy is never the cheapest of the
four methods under the reference cost model. \textbf{We do not claim}
Pareto superiority under ACK-inclusive cost in the Stressed condition: a
pre-registered test of exactly that claim was rejected, and our holdout
data reproduces the rejection with a signed confidence interval. And we do
not claim any hardware result; this is a simulation study, and the
estimator, disturbance and airtime models are stated parameter assumptions
rather than measurements.

The remainder of the paper is organised as follows.
Section~\ref{sec:problem} formalises the setting.
Section~\ref{sec:method} specifies the policy in enough detail to
reimplement, including both algorithms and the receiver-side
acknowledgement semantics. Section~\ref{sec:setup} describes the
pre-registered experimental protocol and the statistical conventions.
Section~\ref{sec:results} reports results, including the design-evolution
ablation and every rejected claim. Section~\ref{sec:discussion} explains
the negative results mechanistically rather than deferring them, and
Section~\ref{sec:limitations} states the boundaries of what has been
shown.
